\documentclass[11pt]{article}
\usepackage[margin=1in,top=0.85in,bottom=0.85in]{geometry}
\usepackage{amsmath,amssymb}
\usepackage[hidelinks]{hyperref}

\DeclareMathOperator{\ord}{ord}
\DeclareMathOperator{\lcm}{lcm}

\title{Two partitions of a prime\protect\\Proof roadmap}
\author{}
\date{}

\begin{document}
\maketitle

The case $q=3$ gives $11=2^3+3=2+3^2$. For $q>3$, the main
reductions are a uniform bound on the odd gap, exclusion of $q\mid M$,
and recovery of the only possible multiplier $h$ from $q$. The remaining
problem is to exclude the resulting candidate as $q$ varies.
The initial square argument and the later use of quadratic reciprocity
belong at different points in the proof.

\section{The integer reduction}

Write $p=2^{c+M}+q^e=2^c+q^{e+N}$ as in the concise notes. The
identities to carry forward are
\[
 Aq^e=2^M-1,\qquad A2^c=q^N-1,\qquad Ap=2^Mq^N-1.
\]
Reduction modulo $3$ gives $2\mid M$. If $N$ is even, write
$Ap=(X-1)(X+1)$ with $X^2=2^Mq^N$. Both factors exceed $A$,
contrary to primality of $p$. This is the work done by the initial
quadratic form. With $N$ odd, $Ap=qX^2-1$ can be read as a signed
norm in $\mathbb Q(\sqrt q)$ and gives $(q/p)=1$. That residue
condition can stay as an auxiliary congruence.

Factoring $2^M-1$ gives $q^e\le 2^{M/2}+1$ and $A\ge 2^{M/2}-1$.
Keep the short arguments for $A=1$, $N=1$, and $q=3$ together.
The rest of the integer reduction gives
\[
 c=v_2(q-1)\text{ even},\quad e\ge2\text{ even},\quad
 \gcd(N,6)=1,\quad N\ge5,\qquad 2^{M/2}<q^N<2^M.
\]
In particular, $3\nmid N$ uses the common divisor argument at $d=3$.
With $r=\ord_q(2)$ and $\rho=v_q(2^r-1)$, lifting gives
$e=\rho+v_q(M)$ and $r\mid(q-1)/3$. Keep the modulo $8$ and
$5$ arguments, which give $q\equiv1,97,193\pmod{240}$
in the remaining branch.

For later use, if a prime $d$ divides both $M$ and $N$, put
$Y=2^{M/d}q^{N/d}$. The size of $p$ gives $\ord_p(Y)=d$
and $Y-1\mid A$. If $d\mid q-1$, the second partition then gives
$1\equiv 2^c+1\pmod d$, a contradiction. Thus with
$E=\lcm(2,r)$ we have $E\mid M$, $E\mid q-1$ and $\gcd(E,N)=1$.
This is the coprimality used below.

\section{Laurent's bound and the finite branch}

The following ordering shortens the bound in the long notes.
First $q^{e+N}>2^{c+M}$, since $q^e>2^c$. Also $q^e\le2^{M/2}+1$,
so
\[
 q^{e+N}>2^{c+M}\ge 2^c(q^e-1)^2>q^{2e},
 \qquad\text{hence } e<N.
\]
Here $c\ge2$ and $q^e>2$. This step applies also when $q\mid M$.

Set $m_1=c+M$, $n_2=e+N$ and $\Lambda=n_2\log q-m_1\log 2$.
The gap equations give $0<\Lambda<(27/26)q^{-N}$.
The starting point was Miyazaki--Scott--Styer's use of this
estimate for exponential equations \cite[Proposition~6.1]{MSS}.
From Laurent's Corollary~2
\cite{Laurent}, take $D=1$, $(m,C_2)=(10,25.2)$ and heights
$H_1=1$, $H_2=\log q$. Then
\[
 b'=\frac{m_1}{\log q}+n_2
 <(1+1/\log 2)n_2<5N.
\]
The factor $\log q$ in Laurent's bound cancels the one in
$N\log q$. The resulting inequality is
\[
 N<25.2\max\{10,\log(5N)+0.38\}^2+1.
\]
At $N=2521$ there is equality in the numerical comparison.
The ratio $(x-1)/\max\{10,\log(5x)+0.38\}^2$ increases
for $x\ge2521$. Thus $N\le2519$ by oddness, $n_2\le5037$,
and $M<5037\log_2 q$.

If $q\mid M$, then $qr\mid M$, and $r>\log_2 q$.
Consequently $q<5037$ and $M<5037\cdot13=65481$.
The existing finite check in \path{scripts/exp-gcd/check_q_divides_gap.lua}
works with these cutoffs. It leaves 31 primes after the initial
conditions, rejects 28 by the minimum $M$, one by $v_3(M)$, and
two modulo $5$. The script with its original, larger limits also
returns zero. This closes the $q\mid M$ branch.

\section{The multiplier determines the candidate}

Now $q\nmid M$, $e=\rho$, and $q\ge97$. Set
\[
 Q=q^\rho,\qquad M=hE,\qquad a=\frac{2^E-1}{Q}.
\]
The full power $Q$ belongs to one of the coprime factors $2^{E/2}\pm1$.
Thus $2^E\ge(Q-1)^2>Q$. From
$q^{n_2}>2^{c+hE}>q^{\rho h}$ we get $h\le1260$.
Since $Q\ge97^2$ and $2h\le2520$, Bernoulli's inequality gives
$(1-1/Q)^{2h}\ge1-2h/Q>1/2$. Hence
\[
 q^{n_2}>2^{c+hE}\ge 2^c(Q-1)^{2h}>Q^{2h},
 \qquad (2h-1)\rho<N,\qquad h\le630.
\]
In particular $h<Q$. These bounds use the single application
of Laurent above.

Since $2^E\equiv1\pmod Q$, the geometric sum gives
\[
 A=a\sum_{j=0}^{h-1}2^{jE}\equiv ah\pmod Q,
 \qquad
 h\equiv-(2^ca)^{-1}\pmod Q.
\]
The last congruence uses $A2^c=q^N-1$ and $N>\rho$.
Thus $h$ is the least positive residue on the right, and a residue
above $630$ excludes $q$. This gives a function of $q$ even before
factoring $a$.

For the candidate $h$ in range define the integer
\[
 F(q)=1+2^c\frac{2^{hE}-1}{Q}.
\]
The remaining equation is $F(q)=q^N$. In particular $N=v_q(F(q))$
and the candidate must satisfy
\[
 5\le v_q(F(q))\le2519,\qquad
 (2h-1)\rho<v_q(F(q)),\qquad \gcd(v_q(F(q)),6E)=1.
\]
These are necessary conditions on the actual integer $F(q)$.
Showing that it always fails one of them would close the proof.
More generally it suffices to rule out $F(q)=q^N$ in this range
under the standing hypotheses. The original problem also requires
$p=2^{c+hE}+Q$ prime. The bound on $N$ reduces the exponent choices
to a finite list. The Roth argument in the long notes proves
finiteness as $q$ varies without an effective cutoff.

\section{Where reciprocity enters}

The full order condition on the part of $A$ fixed by $q$ is
$\ord_a(q)\mid N$, because $a\mid A$ and $q^N\equiv1\pmod A$.
In particular an order sharing a prime factor with $E$, or exceeding
$2519$, excludes that $q$. For a surviving $h$, the valuation above
also fixes the exponent to use in this condition.

\paragraph{Quadratic reciprocity.}
The supplement for $2$ already enters in the modulo $8$ reduction.
For each prime $\ell\mid A$, the order of $q$ modulo $\ell$ is odd,
so $(q/\ell)=1$. Since $q\equiv1\pmod4$, we get $(\ell/q)=1$.
Thus every positive divisor of $A$ is a square modulo $q$. Taking
$a$ and using $2^cah\equiv-1\pmod q$ shows that $h$ is a square.
This is a short lemma worth keeping. A quadratic field also occurs
through the unit in $\mathbb Q(\sqrt5)$ in the conductor-five test.
That belongs with the cyclotomic argument.

\paragraph{Cyclotomic fields and power reciprocity.}
Pomerance \cite[\S2]{PomCycl} gives the starting point for studying
the prime divisors of $\Phi_d(2)$ by their orders at base $2$.
Brent \cite[\S\S2--3]{Brent} develops the Gauss and Aurifeuillean
factorizations used in the longer notes.
For a prime power $m=s^b\mid r$, put $B=2^{E/m}$.
The norm of $B-\zeta_m$ selects $\Phi_m(B)$, whose prime factors
away from $q,s$ divide $a$. The power-residue symbol at each such
prime ideal must be $1$ by the order condition. The reciprocity
product computes their product from local terms above $q$
after the terms at $s$ have been shown trivial. A nontrivial product
proves an order obstruction without factoring $a$. The concise
notes give the explicit formula, adapted from Sharifi's use of
reciprocity on cyclotomic norms \cite[\S3, (5)]{SharifiCycl}.
Equality to $1$ can come from cancellation or from the exponent
of the symbol failing to detect the full $s$-part of an order.

The test can also prescribe a power class for $h$.
At conductor five, the separate factor gives the Fibonacci
condition $q\mid F_{(q-1)/5}$ (here $F_j$ is the Fibonacci
sequence), and the grouped test then requires $h$ to be a fifth
power modulo $q$. When the Fibonacci condition holds, $h=1$
passes both. A general reciprocity proof still needs to show
that some product is nontrivial, or that the primes passing the
products fail a further necessary condition such as the required
valuation of $F(q)$.

\section{Working proof and Lean}

\begin{enumerate}
\item Formalize the integer reduction through $e<N$, including the
$q=3$ case. Use the ordered partitions as the input and keep the
three equations for $A$ together.
\item State the specialization of Laurent as an explicit dependency.
Prove the integer consequences of the bound and replace the Lua
check by a certified finite calculation if using Lean for the
whole argument. The smaller check needs only the valuation at $3$
and the order test at $5$ after the size test.
\item Prove the multiplier bound, inverse congruence, and reduction
to $F(q)=q^N$. This is the endpoint of the current proof.
\item Keep the elementary quadratic lemma available. Add the
general cyclotomic reciprocity calculation to the formal proof
when a final argument uses it.
\end{enumerate}

For a shorter natural-language proof, retain the common-divisor
lemmas needed for $3\nmid N$ and $\gcd(E,N)=1$. Keep the wider
coprimality casework, forbidden odd divisors, pair scans, conductor
examples and alternative approximation bounds in the research notes.
Their main role so far is to exclude families or to test candidate
arguments for the last step.

\bibliographystyle{alpha}
\bibliography{prime_power_notes}
\end{document}
